\documentclass[10pt,a4paper]{amsart}
\usepackage{amsmath,amssymb,amsthm,graphicx}
% Retained minimal preamble for the Stacks Project chapter more-algebra.tex (More on Algebra).
% The chapter file pinned-original-source.tex carries only \input{preamble} on its first line, so no
% preamble text lives inside the chapter itself.  The Stacks Project uses one single shared preamble for
% every chapter; the snapshot stages it identically as sources/stacks-chow/preamble.tex and
% sources/stacks-etale-cohomology/preamble.tex.  The block below is copied mechanically, line for line,
% from sources/stacks-chow/preamble.tex lines 14-18 (the Xy-pic setup) and 154-217 (the native theorem
% declarations together with the mathematical macros).  Nothing was added, removed or redefined; the
% wrapper class and the amsmath/amssymb/amsthm/graphicx base packages are supplied by the assembler.
% Chapter: more-algebra.tex -- More on Algebra.
% For commutative diagrams we use Xy-pic
\usepackage[all]{xy}

% We use 2cell for 2-commutative diagrams.
\xyoption{2cell}

% Theorem environments.
%
\theoremstyle{plain}
\newtheorem{theorem}[subsection]{Theorem}
\newtheorem{proposition}[subsection]{Proposition}
\newtheorem{lemma}[subsection]{Lemma}

\theoremstyle{definition}
\newtheorem{definition}[subsection]{Definition}
\newtheorem{example}[subsection]{Example}
\newtheorem{exercise}[subsection]{Exercise}
\newtheorem{situation}[subsection]{Situation}

\theoremstyle{remark}
\newtheorem{remark}[subsection]{Remark}
\newtheorem{remarks}[subsection]{Remarks}

\numberwithin{equation}{subsection}

% Macros
%
\def\lim{\mathop{\mathrm{lim}}\nolimits}
\def\colim{\mathop{\mathrm{colim}}\nolimits}
\def\Spec{\mathop{\mathrm{Spec}}}
\def\Hom{\mathop{\mathrm{Hom}}\nolimits}
\def\Ext{\mathop{\mathrm{Ext}}\nolimits}
\def\SheafHom{\mathop{\mathcal{H}\!\mathit{om}}\nolimits}
\def\SheafExt{\mathop{\mathcal{E}\!\mathit{xt}}\nolimits}
\def\Sch{\mathit{Sch}}
\def\Mor{\mathop{\mathrm{Mor}}\nolimits}
\def\Ob{\mathop{\mathrm{Ob}}\nolimits}
\def\Sh{\mathop{\mathit{Sh}}\nolimits}
\def\NL{\mathop{N\!L}\nolimits}
\def\CH{\mathop{\mathrm{CH}}\nolimits}
\def\proetale{{pro\text{-}\acute{e}tale}}
\def\etale{{\acute{e}tale}}
\def\QCoh{\mathit{QCoh}}
\def\Ker{\mathop{\mathrm{Ker}}}
\def\Im{\mathop{\mathrm{Im}}}
\def\Coker{\mathop{\mathrm{Coker}}}
\def\Coim{\mathop{\mathrm{Coim}}}

% Boxtimes
%
\DeclareMathSymbol{\boxtimes}{\mathbin}{AMSa}{"02}

%
% Macros for moduli stacks/spaces
%
\def\QCohstack{\mathcal{QC}\!\mathit{oh}}
\def\Cohstack{\mathcal{C}\!\mathit{oh}}
\def\Spacesstack{\mathcal{S}\!\mathit{paces}}
\def\Quotfunctor{\mathrm{Quot}}
\def\Hilbfunctor{\mathrm{Hilb}}
\def\Curvesstack{\mathcal{C}\!\mathit{urves}}
\def\Polarizedstack{\mathcal{P}\!\mathit{olarized}}
\def\Complexesstack{\mathcal{C}\!\mathit{omplexes}}
% \Pic is the operator that assigns to X its picard group, usage \Pic(X)
% \Picardstack_{X/B} denotes the Picard stack of X over B
% \Picardfunctor_{X/B} denotes the Picard functor of X over B
\def\Pic{\mathop{\mathrm{Pic}}\nolimits}
\def\Picardstack{\mathcal{P}\!\mathit{ic}}
\def\Picardfunctor{\mathrm{Pic}}
\def\Deformationcategory{\mathcal{D}\!\mathit{ef}}

% --- appended 2026-10-12 (Chief): the project own definitions of the reference/slogan/history
% display environments, copied mechanically from the shared Stacks preamble lines 8-12
% (\usepackage{verbatim} plus the three \newenvironment lines, which the chapter preamble keeps
% verbatim from the project preamble). Nothing added, removed or redefined: these environments
% carry no mathematical content and the project itself typesets them as comments. Without these
% three lines a counted statement whose own text uses them fails with Environment slogan undefined.
\usepackage{verbatim}
\newenvironment{reference}{\comment}{\endcomment}
\newenvironment{slogan}{\comment}{\endcomment}
\newenvironment{history}{\comment}{\endcomment}
\usepackage{cleveref}
\crefname{originallemma}{original lemma}{original lemmas}
\Crefname{originallemma}{Original lemma}{Original lemmas}
\makeatletter
\newlabel{fields-section-Kummer}{{09I6}{source}{}{sourcepointer-1}{}}
\newlabel{fields-section-Kummer@cref}{{[originallemma][0][]09I6}{[1][0][]source}{}{}{}}
\makeatother
\crefname{originallemma}{original lemma}{original lemmas}
\Crefname{originallemma}{Original lemma}{Original lemmas}
\makeatletter
\newlabel{lemma-prepare}{{34977}{source}{}{sourcepointer-2}{}}
\newlabel{lemma-prepare@cref}{{[originallemma][0][]34977}{[1][0][]source}{}{}{}}
\makeatother
\crefname{originalsection}{original section}{original sections}
\Crefname{originalsection}{Original section}{Original sections}
\makeatletter
\newlabel{section-ramification}{{33172}{source}{}{sourcepointer-3}{}}
\newlabel{section-ramification@cref}{{[originalsection][0][]33172}{[1][0][]source}{}{}{}}
\makeatother
\title{An extension obtained by adjoining a root of the prepared polynomial is Galois and falls into one of four cases according to the valuation of the parameter}
\author{The Stacks Project Authors}
\date{stacks-project:more-algebra.tex:local-snapshot-20261012; GNU FDL 1.2 or later}
\begin{document}
\maketitle
\setcounter{subsection}{1}
\noindent\textit{Editorial note.} All delivered mathematical text is the author's own, copied from the pinned Stacks Project chapter file more-algebra.tex; no statement or proof step was rewritten, reordered or removed.\par\medskip
\noindent\textit{Editorial note.} The chapter file pinned-original-source.tex carries only \textbackslash{}input\{preamble\} on its first line, so no preamble text lives inside the chapter. The retained file sources/stacks-more-algebra/retained-preamble.tex is a mechanical, line-for-line copy of the project's single shared preamble as staged in the snapshot; nothing was added, removed or redefined and no package was added.\par\medskip
\noindent\textit{Editorial note.} The source URL is the source-level Stacks Project URL, as required when more than one card is produced from one source; the per-claim tag is carried by claim\_key.\par\medskip
\noindent\textit{Editorial note.} The author's label for this statement contains the misspelling 'polynial' for 'polynomial'; it is reproduced unchanged in claim\_key, and the card's own prose uses the correct spelling. The three references inside the counted proof span are recorded as non-load-bearing navigation pointers.\par\medskip
\noindent\textit{Allowed standard prerequisites (editorial).}\begin{itemize}
\item Kummer theory in the presence of a primitive p-th root of unity. invoked at proof lines 35045-35048. Gives that the extension is Galois of order dividing p, which makes the four-case analysis exhaustive.
\item The preparation lemma's polynomial P and its addition formula. used at statement line 35020 and proof lines 35045-35046. Supplies the polynomial whose root is adjoined and the rewriting of that adjunction.
\item Ramification theory for extensions of a discrete valuation ring, including the multiplicativity of ramification indices and the description of weakly unramified extension. used at proof lines 35049-35051 and 35058-35071. Is the framework in which the four cases are separated and the implications (A), (B) and (C) are established.
\end{itemize}\medskip
\ifdefined\hypertarget\hypertarget{sourcepointer-1}{}\fi
\noindent\textit{Source reference.} Original lemma 09I6: focused-seeds/AGP/sources/stacks-snapshot-20261011/tags.txt, source lines 12029--12029. Cross-chapter navigation pointer into the Stacks Project; the standard fact it supplies is applied inside the counted proof span and the argument resting on it stays there.\par\medskip
\ifdefined\hypertarget\hypertarget{sourcepointer-2}{}\fi
\noindent\textit{Source reference.} Original lemma 34977: sources/stacks-more-algebra/pinned-original-source.tex, source lines 34977--34977. Same-chapter navigation pointer; the statement it names is standard for the section and the argument that consumes it stays inside the counted proof span.\par\medskip
\ifdefined\hypertarget\hypertarget{sourcepointer-3}{}\fi
\noindent\textit{Source reference.} Original section 33172: sources/stacks-more-algebra/pinned-original-source.tex, source lines 33172--33172. Same-chapter navigation pointer; the statement it names is standard for the section and the argument that consumes it stays inside the counted proof span.\par\medskip
\begin{lemma}
\label{lemma-extension-defined-by-nice-polynial}
Let $A$ be a discrete valuation ring of mixed characteristic $(0, p)$
which contains a primitive $p$th root of $1$.
Let $P(t) \in A[t]$ be the polynomial of Lemma \ref{lemma-prepare}.
Let $\xi \in K$.
Let $L$ be an extension of $K$ obtained by
adjoining a root of $P(z) = \xi$. Then $L/K$ is Galois and one of the
following happens
\begin{enumerate}
\item $L = K$,
\item $L/K$ is unramified with respect to $A$ of degree $p$,
\item $L/K$ is totally ramified with respect to $A$
with ramification index $p$, and
\item the integral closure $B$ of $A$ in $L$ is a discrete valuation ring,
$A \subset B$ is weakly unramified, and $A \to B$ induces a purely inseparable
residue field extension of degree $p$.
\end{enumerate}
Let $\pi$ be a uniformizer of $A$. We have the following implications:
\begin{enumerate}
\item[(A)] If $\xi \in A$, then we are in case (1) or (2).
\item[(B)] If $\xi = \pi^{-n}a$ where $n > 0$ is not divisible by
$p$ and $a$ is a unit in $A$, then we are in case (3)
\item[(C)] If $\xi = \pi^{-n} a$ where $n > 0$ is divisible by $p$ and
the image of $a$ in $\kappa_A$ is not a $p$th power, then we are in case (4).
\end{enumerate}
\end{lemma}

\begin{proof}
Adjoining a root of $P(z) = \xi$ is the same thing as adjoining a root
of $y^p = w^p(1 + \xi)$. Since $K$ contains a primitive $p$th root of $1$
the extension is Galois of order dividing $p$ by the discussion in
Fields, Section \ref{fields-section-Kummer}.
It immediately follows from the discussion in
Section \ref{section-ramification} that we are in one of the cases (1) -- (4)
listed in the lemma.

\medskip\noindent
Case (A). Here we see that $A \to A[x]/(P(x) - \xi)$ is a finite
\'etale ring extension. Hence we are in cases (1) or (2).

\medskip\noindent
Case (B). Write $\xi = \pi^{-n}a$ where $p$ does not divide $n$.
Let $B \subset L$ be the integral closure of $A$ in $L$.
If $C = B_\mathfrak m$ for some maximal ideal $\mathfrak m$,
then it is clear that $p \text{ord}_C(z) = -n \text{ord}_C(\pi)$.
In particular $A \subset C$ has ramification index divisible by $p$.
It follows that it is $p$ and that $B = C$.

\medskip\noindent
Case (C). Set $k = n/p$. Then we can rewrite the equation as
$$
(\pi^kz)^p - \pi^{n - k} (\pi^kz) + \sum a_i \pi^{n - ik} (\pi^kz)^i = a
$$
Since $A[y]/(y^p - \pi^{n - k}y - \sum  a_i \pi^{n - ik} y^i - a)$
is a discrete valuation ring weakly unramified over $A$, the lemma follows.
\end{proof}

\end{document}
